\documentclass[a4paper,12pt]{article}
\usepackage[T2A]{fontenc}
\usepackage[utf8]{inputenc}
\usepackage[russian]{babel}
\usepackage{amsmath,amssymb,mathtools}
\usepackage{PTSans}
\renewcommand{\familydefault}{\sfdefault}
\usepackage{icomma}
\usepackage[left=23mm,right=23mm,top=23mm,bottom=23mm,headheight=25pt,headsep=6mm,footskip=12mm]{geometry}
\usepackage{microtype}
\usepackage{tikz}
\usetikzlibrary{arrows.meta,positioning,shapes.geometric,calc}
\usepackage{tabularx,booktabs}
\usepackage{enumitem}
\usepackage{fancyhdr}
\usepackage{titlesec}
\usepackage{setspace}
\usepackage{needspace}
\usepackage{hyperref}
\usepackage{xcolor}
\definecolor{accent}{HTML}{008080}
\definecolor{darkaccent}{HTML}{00575B}
\definecolor{textgray}{HTML}{343A3B}
\definecolor{muted}{HTML}{667579}
\hypersetup{hidelinks,pdftitle={Чек-лист: отношения, доли и пропорции},pdfauthor={Лёвин Артём Александрович}}
\setstretch{1.18}
\setlength{\parindent}{0pt}
\setlength{\parskip}{6pt plus 1pt minus 1pt}
\setlength{\headheight}{25pt}
\emergencystretch=1.5em
\setlist[enumerate]{leftmargin=*,itemsep=5pt,topsep=5pt}
\setlist[itemize]{leftmargin=*,itemsep=3pt,topsep=3pt}
\titleformat{\section}{\large\bfseries\color{darkaccent}}{\thesection.}{0.6em}{}[\vspace{2pt}\color{accent}\titlerule]
\titlespacing*{\section}{0pt}{16pt}{11pt}
\titleformat{\subsection}{\normalsize\bfseries\color{darkaccent}}{}{0pt}{}
\titlespacing*{\subsection}{0pt}{12pt}{3pt}
\pagestyle{fancy}
\fancyhf{}
\fancyhead[L]{\small\color{textgray}Отношения, доли и пропорции}
\fancyhead[R]{\small\color{textgray}Ксения Васильченко}
\fancyfoot[C]{\small\thepage}
\renewcommand{\headrulewidth}{0.4pt}
\renewcommand{\footrulewidth}{0pt}
\makeatletter
\renewcommand{\headrule}{\hbox to\headwidth{\color{accent}\leaders\hrule height \headrulewidth\hfill}\vskip-\headrulewidth}
\makeatother
\newcommand{\keyidea}[1]{\textcolor{darkaccent}{\textbf{Главная мысль.}}\ #1}
\newcommand{\ex}[1]{\par\textcolor{darkaccent}{\textbf{Пример.}}\ #1}
\newcommand{\warn}[1]{\textcolor{darkaccent}{\textbf{Проверьте.}}\ #1}
\begin{document}
\begin{titlepage}
\thispagestyle{empty}
\centering
\vspace*{27mm}
{\large\color{darkaccent}МАТЕМАТИКА\par}
\vspace{20mm}
{\LARGE\color{textgray}Чек-лист\par}
\vspace{13mm}
{\LARGE\color{accent}\bfseries Отношения, доли\par и пропорции\par}
\vspace{20mm}
{\normalsize\color{textgray}Как выбирать целое, сравнивать величины\par и находить проценты\par}
\vfill
{\large\color{textgray}Ксения Васильченко\par}
\vspace{10mm}
{\normalsize\color{textgray}17 ноября 2026 г.\par}
\vspace{14mm}
{\small\color{textgray}Лёвин Артём Александрович\par эксклюзивно для Ксении Васильченко\par}
\vspace{14mm}
\begin{tikzpicture}[x=1cm,y=1cm]
 \draw[accent,line width=.85pt] (0,0) .. controls (1.1,.33) and (2.0,-.30) .. (3.1,0)
  .. controls (4.2,.36) and (5.2,-.31) .. (6.35,0)
  .. controls (7.2,.25) and (8.0,-.17) .. (9.0,0);
\end{tikzpicture}
\vspace*{2mm}
\end{titlepage}

\section{Что показывает отношение}

\keyidea{Чтобы найти \emph{отношение величины $A$ к величине $B$}, разделите первую величину на вторую:}
\[
\boxed{\displaystyle A:B=\frac{A}{B}},\qquad B\ne0.
\]
\textbf{Порядок важен:} отношение $A$ к $B$ и отношение $B$ к $A$ обычно различаются.

Если речь идёт о \emph{части целого}, в числителе записывают размер части, а в знаменателе --- размер всего целого. Для неотрицательных величин такая доля лежит между нулём и единицей.

\ex{Решено $4$ задачи из $5$. Доля решённых задач равна}
\[
\frac{4}{5}=0{,}8=80\%.
\]
Запись $4$ сама по себе не показывает, много это или мало: требуется знать общее количество задач.

\textbf{Три формулировки --- один приём деления:}
\begin{itemize}
\item «Какую часть $B$ составляет $A$?» $\longrightarrow\ \dfrac{A}{B}$ (если $A$ --- часть $B$).
\item «Найдите отношение $A$ к $B$» $\longrightarrow\ \dfrac{A}{B}$.
\item «Во сколько раз $A$ больше $B$?» $\longrightarrow\ \dfrac{A}{B}$, если $A>B>0$.
\end{itemize}
Число $\frac{A}{B}>1$ допустимо для отношения и ответа «во сколько раз», но \emph{не является долей части собственного целого}.

\section{Как правильно выбрать целое}

Прежде чем делить, установите, \emph{от какой совокупности} требуется найти долю. Иногда её размер не дан напрямую --- сначала нужно сложить части.

\ex{Из $480\,\text{м}^2$ катка расчищено $320\,\text{м}^2$. Расчищенная доля:}
\[
\frac{320}{480}=\frac{2}{3}.
\]
Знаменатель равен $480$, поскольку спрашивают часть \emph{всего катка}.

\ex{В букете $24$ красных и $18$ белых тюльпанов. Всего $24+18=42$ тюльпана. Поэтому}
\[
\text{доля красных}=\frac{24}{42}=\frac47,\qquad
\text{доля белых}=\frac{18}{42}=\frac37.
\]
Сумма долей равна $\frac47+\frac37=1$: учтены все цветы.

\ex{Сплав содержит $3{,}76$ кг олова и $0{,}8$ кг сурьмы. Масса всего сплава $4{,}56$ кг. Доля олова:}
\[
\frac{3{,}76}{3{,}76+0{,}8}=\frac{3{,}76}{4{,}56}=\frac{47}{57}.
\]
Деление $3{,}76:0{,}8$ дало бы \emph{отношение олова к сурьме}, а не долю олова во всём сплаве.

\section{«Во сколько раз» и обратное отношение}

\keyidea{Для вопроса «во сколько раз $A$ больше $B$?» делите $A$ на $B$. Не меняйте числа местами.}

\ex{В задаче о катке весь каток больше расчищенной площади в}
\[
\frac{480}{320}=\frac32=1{,}5\ \text{раза}.
\]
Сравните результаты: доля очищенного катка --- $\frac23$, а отношение всего катка к очищенной части --- $\frac32$. Это взаимно обратные числа.

\ex{Красных тюльпанов $24$, белых $18$. Отношение красных к белым равно}
\[
\frac{24}{18}=\frac43=1\frac13.
\]
Оно больше единицы, поскольку красных цветов больше, чем белых. Это \emph{не} доля красных цветов во всём букете.

\section{Сначала приведите величины к одним единицам}

\warn{Сравниваемые однородные величины должны быть выражены в одинаковых единицах.}
\[
1\,\text{кг}=1000\,\text{г},\qquad 4{,}4\,\text{кг}=4400\,\text{г}.
\]
\ex{Котёнок весит $99$ г, кошка --- $4{,}4$ кг. Отношение масс котёнка к кошке:}
\[
\frac{99\,\text{г}}{4400\,\text{г}}=\frac{99}{4400}=\frac9{400}.
\]
Единицы сократились; отношение получилось безразмерным.

При умножении килограммов на $1000$ десятичная запятая смещается на три разряда вправо:
\[
0{,}75\,\text{кг}=750\,\text{г},\qquad
1{,}8\,\text{кг}=1800\,\text{г},\qquad
13{,}125\,\text{кг}=13125\,\text{г}.
\]

\Needspace{7\baselineskip}
\section{Как найти остаток целого}

Если всё целое принято за $1$, то оставшаяся доля равна $1$ минус известная доля.

\ex{Первая бригада прошла $\frac59$ тоннеля, вторая --- остальную часть:}
\[
\text{вторая бригада:}\quad 1-\frac59=\frac49.
\]
Участок первой бригады больше участка второй в
\[
\frac{\frac59}{\frac49}=\frac59\cdot\frac94=\frac54=1{,}25\ \text{раза}.
\]
\textbf{Почему именно так?} Сначала нашли обе сравниваемые части одного целого; затем разделили долю первой бригады на долю второй.

\section{Равные отношения и масштабирование рецепта}

\textbf{Пропорция} --- равенство двух отношений:
\[
\frac{a}{b}=\frac{c}{d},\qquad b\ne0,\ d\ne0.
\]
Если обе сравниваемые величины умножить на одно и то же положительное число, их отношение не изменится:
\[
\frac{ka}{kb}=\frac ab,\qquad k>0,\ b\ne0.
\]
\ex{Для блюда используют $0{,}7$ кг моркови и $4{,}9$ кг капусты. Отношение масс:}
\[
0{,}7:4{,}9=1:7.
\]
На каждый $1$ кг моркови нужно $7$ кг капусты. Для $2$ кг моркови потребуется $2\cdot7=14$ кг капусты. Пропорции ингредиентов сохранились.

\section{Как сравнить результаты при разных объёмах}

Когда количество попыток отличается, сравнивают не число успехов, а \emph{долю успешных попыток}.

\ex{У Кати $15$ удачных прыжков из $18$, у Даши --- $19$ из $21$.}
\[
\text{Катя: }\frac{15}{18}=\frac56=\frac{35}{42},\qquad
\text{Даша: }\frac{19}{21}=\frac{38}{42}.
\]
Поскольку $\frac{38}{42}>\frac{35}{42}$, доля удачных прыжков у Даши выше. Сравнивать только числа $19$ и $15$ нельзя: количество попыток различается.

\section{Проценты от нужной группы}

\keyidea{Если ответ просят в процентах, сначала найдите \emph{отношение к указанной группе}, а затем умножьте его на $100\%$.}
\[
\boxed{\displaystyle \text{Искомый процент}=\frac{A}{B}\cdot100\%}.
\]
\ex{Футболом занимаются $44\%$ ребят, баскетболом --- $27\%$, лёгкой атлетикой --- остальные $29\%$.}

Среди занимающихся \emph{игровыми видами спорта} футболисты составляют
\[
\frac{44}{44+27}\cdot100\%=\frac{44}{71}\cdot100\%\approx61{,}97\%.
\]
Среди занимающихся \emph{футболом или лёгкой атлетикой} футболисты составляют
\[
\frac{44}{44+29}\cdot100\%=\frac{44}{73}\cdot100\%\approx60{,}27\%.
\]
Знаменатели разные, потому что в каждом вопросе рассматривается своя группа.

\section{Типичные ошибки и быстрый контроль}

\renewcommand{\arraystretch}{1.24}
\begin{tabularx}{\linewidth}{@{}p{.38\linewidth}X@{}}
\toprule
\textbf{Ошибка} & \textbf{Что сделать вместо этого}\\
\midrule
Перепутаны числитель и знаменатель & Подчеркните слова «$A$ к $B$»: делим именно $A$ на $B$.\\
Часть поделена на другую часть & Проверьте, не требуется ли в знаменателе всё целое.\\
Килограммы делят на граммы & Сначала переведите данные в одинаковые единицы.\\
Сравнивают число успехов & Сравните долю успехов от числа всех попыток.\\
Забыли перевести ответ в проценты & Умножьте найденное отношение на $100\%$.\\
\bottomrule
\end{tabularx}

\textbf{Контроль смысла:} доля части целого не больше $1$; обратное отношение «во сколько раз больше» больше $1$; сумма долей всех непересекающихся частей равна $1$.

\clearpage
\thispagestyle{empty}
\begin{center}
 {\Large\bfseries\color{darkaccent}Алгоритм: как решить задачу на отношение}
\end{center}
\vspace{7mm}
\begin{center}
\tikzset{
 startstop/.style={ellipse,draw=darkaccent,thick,minimum width=48mm,minimum height=10mm,align=center,font=\normalsize},
 action/.style={rectangle,rounded corners=2pt,draw=darkaccent,thin,text width=93mm,minimum height=11mm,align=center,inner sep=3pt,font=\normalsize},
 smallaction/.style={rectangle,rounded corners=2pt,draw=darkaccent,thin,text width=43mm,minimum height=13mm,align=center,inner sep=3pt,font=\small},
 decision/.style={diamond,aspect=2.4,draw=accent,thick,text width=42mm,align=center,inner sep=1.5pt,font=\small},
 arrow/.style={-{Stealth[length=3mm]},thin,draw=darkaccent}
}
\begin{tikzpicture}[node distance=7mm and 13mm]
 \node[startstop] (start) {Прочитайте условие};
 \node[action,below=of start] (choose) {Назовите $A$: что сравниваем?\\ Назовите $B$: с чем сравниваем?};
 \node[decision,below=of choose,yshift=-1mm] (group) {Размер $B$ известен?};
 \node[smallaction,below left=9mm and 9mm of group] (sum) {Нет: найдите $B$\\ (например, сумму частей)};
 \node[smallaction,below right=9mm and 9mm of group] (known) {Да: используйте\\ известное $B$};
 \coordinate (merge) at ($(sum.south)!0.5!(known.south)+(0,-5mm)$);
 \node[action,below=7mm of merge] (units) {Приведите однородные величины к одинаковым единицам. Убедитесь, что $B\ne0$.};
 \node[action,below=of units] (divide) {Вычислите $A:B=\dfrac AB$ и сократите дробь.};
 \node[decision,below=of divide,yshift=-1mm] (percent) {Ответ нужен в процентах?};
 \node[smallaction,below left=9mm and 9mm of percent] (yes) {Да: умножьте\\ на $100\%$};
 \node[smallaction,below right=9mm and 9mm of percent] (no) {Нет: оставьте\\ отношение};
 \coordinate (merge2) at ($(yes.south)!0.5!(no.south)+(0,-5mm)$);
 \node[startstop,below=5mm of merge2] (done) {Проверьте смысл ответа};
 \draw[arrow] (start) -- (choose);
 \draw[arrow] (choose) -- (group);
 \draw[arrow] (group.south) |- (sum.north) node[pos=.80,above left,font=\small] {нет};
 \draw[arrow] (group.south) |- (known.north) node[pos=.80,above right,font=\small] {да};
 \draw[arrow] (sum.south) |- (merge);
 \draw[arrow] (known.south) |- (merge);
 \draw[arrow] (merge) -- (units);
 \draw[arrow] (units) -- (divide);
 \draw[arrow] (divide) -- (percent);
 \draw[arrow] (percent.south) |- (yes.north) node[pos=.8,above left,font=\small] {да};
 \draw[arrow] (percent.south) |- (no.north) node[pos=.8,above right,font=\small] {нет};
 \draw[arrow] (yes.south) |- (merge2);
 \draw[arrow] (no.south) |- (merge2);
 \draw[arrow] (merge2) -- (done);
\end{tikzpicture}
\end{center}
\clearpage

\section*{Тренировочный блок}
\textcolor{muted}{\small Путь А. Десять заданий: от прямого отношения к сравнению долей и пропорциям. Все ответы записывайте в сокращённом виде.}

\subsection*{Базовый уровень}
\begin{enumerate}[label=\textbf{\arabic*.},itemsep=8pt]
\item Из $12$ задач решены $9$. Какую часть всех задач составляют решённые?
\item Из $540\,\text{м}^2$ катка расчищено $360\,\text{м}^2$. Найдите: а) расчищенную долю; б) во сколько раз весь каток больше расчищенной части.
\item В букете $30$ красных и $20$ белых тюльпанов. Найдите: а) отношение числа красных к числу белых; б) долю белых среди всех тюльпанов.
\end{enumerate}
\subsection*{Средний уровень}
\begin{enumerate}[label=\textbf{\arabic*.},start=4,itemsep=8pt]
\item Сплав состоит из $2{,}4$ кг олова и $0{,}6$ кг другого металла. Какую часть массы сплава составляет олово?
\item Масса котёнка --- $120$ г, кошки --- $4{,}8$ кг. Найдите отношение массы котёнка к массе кошки.
\item Первая бригада построила $\frac7{12}$ тоннеля, вторая --- остальное. Во сколько раз участок первой бригады длиннее участка второй?
\item Для приготовления блюда смешивают $0{,}8$ кг моркови и $5{,}6$ кг капусты. Найдите: а) отношение моркови к капусте в виде $1:n$; б) массу капусты для $2{,}4$ кг моркови при той же пропорции.
\end{enumerate}
\subsection*{Уровень выше среднего}
\begin{enumerate}[label=\textbf{\arabic*.},start=8,itemsep=8pt]
\item Лена выполнила успешно $14$ попыток из $20$, а Маша --- $18$ из $24$. У кого доля удачных попыток больше?
\item В спортивной секции $36\%$ ребят занимаются футболом, $24\%$ --- баскетболом, остальные --- лёгкой атлетикой. Какой процент \emph{среди занимающихся игровыми видами спорта} составляют футболисты?
\item В группе А успешны $18$ попыток из $24$, в группе Б --- $20$ из $25$. Определите: а) в какой группе выше доля успехов; б) во сколько раз большая доля превышает меньшую.
\end{enumerate}

\clearpage
\section*{Ответы к тренировочному блоку}
\vspace{2mm}
\renewcommand{\arraystretch}{1.44}
\begin{tabularx}{\linewidth}{@{}p{.12\linewidth}X@{}}
\toprule
\textbf{№} & \textbf{Ответ}\\
\midrule
1 & $\dfrac34$.\\
2 & а) $\dfrac23$; б) $\dfrac32=1{,}5$ раза.\\
3 & а) $\dfrac32$ (или $3:2$); б) $\dfrac25$.\\
4 & $\dfrac45$.\\
5 & $\dfrac1{40}$.\\
6 & $\dfrac75=1{,}4$ раза.\\
7 & а) $1:7$; б) $16{,}8$ кг.\\
8 & У Маши: $\dfrac34>\dfrac7{10}$.\\
9 & $60\%$.\\
10 & а) В группе Б: $\dfrac45>\dfrac34$; б) $\dfrac{16}{15}$ раза.\\
\bottomrule
\end{tabularx}

\vspace{13mm}
{\large\bfseries\color{darkaccent}Итоговая памятка}\par
\vspace{2mm}\color{accent}\rule{\linewidth}{0.45pt}\color{black}

\begin{itemize}[label=$\square$,itemsep=7pt]
\item Я понимаю, что именно требуется сравнить: $A$ с $B$.
\item Для доли я выбрал(а) в знаменателе всё указанное целое.
\item Перед делением я привёл(а) однородные величины к общим единицам.
\item Если спрашивают «во сколько раз», я сохранил(а) порядок деления.
\item Если ответ требуется в процентах, я умножил(а) отношение на $100\%$.
\item Я проверил(а), согласуется ли ответ со смыслом задачи.
\end{itemize}
\end{document}
